\section{Limitations and Future Work}
While AQUA-SENTINEL V9 achieves high operational precision, several technical boundary conditions remain for future research.

\subsection{Environmental Constraints and Sea-State Saturation}
The primary technical limitation is observed during extreme sea-state saturation ($>18$ m/s wind speeds). \textit{Rationale:} Under intense surface agitation, the backscatter intensity gradients generated by mechanical spills are physically masked by the white-cap induced speckle noise. In these scenarios, the Laplacian derivative mapping becomes saturated, causing the auxiliary head to misidentify valid signals as stochastic oceanic lookalikes. 

\subsection{Modal and Temporal Resolution Gap}
Current validation is restricted to single-acquisition SAR images. \textit{Rationale:} While the 14ms inference runtime is ideal for real-time edge alerts, the model lacks temporal awareness for resolving the evolution of biogenic films vs. petroleum slicks. Future iterations will explore the integration of multi-temporal SAR sequences to leverage the distinct temporal persistence characteristics of mechanical crude oil as a suppressive modality.

\subsection{Dependency on Ground-Truth Quality}
The 99.66\% precision is inherently coupled to the quality of the multi-modal segmentations used for curriculum training. \textit{Rationale:} Errors in the manual annotation of "lookalike" biogenic films are transmitted into the auxiliary suppression head's latent weights. Subsequent work will incorporate semi-supervised contrastive learning to reduce the system's reliance on high-precision manual labeling in geographically remote oceanic corridors.
